% ============================================================
% Preambolo comune della dispensa del Laboratorio di RO
% ============================================================
\usepackage[utf8]{inputenc}
\usepackage[T1]{fontenc}
\usepackage[english]{babel}
\emergencystretch=3em  % evita overfull nella prosa inglese
\usepackage[a4paper,top=2.6cm,bottom=3.4cm,left=2.6cm,right=2.6cm,footskip=1.7cm]{geometry}
\usepackage{amsmath,amssymb,amsthm,mathtools}
\usepackage{booktabs}
\usepackage{graphicx}
% figures are also reachable from solutions/ (shared preamble)
\graphicspath{{./}{figure/}}
\usepackage{xcolor}
\usepackage{enumitem}
\usepackage{caption}
\usepackage[most]{tcolorbox}
\usepackage{listings}
\usepackage{etoolbox}
\usepackage{titlesec}
\usepackage{pgfplots}
\pgfplotsset{compat=1.16}
\usetikzlibrary{shapes.geometric,arrows.meta}
\usepackage{hyperref}
\hypersetup{
  colorlinks=true,
  linkcolor=blunotte,
  citecolor=blunotte,
  urlcolor=teal,
  filecolor=teal,
  bookmarks=true, bookmarksopen=true, bookmarksopenlevel=1, bookmarksnumbered=true,
  pdftitle={MIP Modelling},
  pdfauthor={Fabio Furini},
  pdfsubject={Mixed-integer linear models for Management Engineering},
  pdfkeywords={operations research, integer linear programming, modelling, binary variables, Gurobi, Python},
  pdfdisplaydoctitle=true
}

% ---------- intestazione e piè di pagina (style of Fabio Furini's lecture notes) ----------
% il logo: chi importa il preambolo da un'altra cartella ridefinisce il percorso
\providecommand{\logosapienza}{figure/sapienza.jpeg}
\usepackage{fancyhdr}
\pagestyle{fancy}
\renewcommand{\footrulewidth}{0.4pt}
\fancyhead{}
\fancyhead[RO,LE]{\href{https://github.com/fabiofurini/mip-modelling}{\textbf{MIP Modelling}}}
\fancyfoot{}
\fancyfoot[LE,RO]{\thepage}
\fancyfoot[LO,CE]{\textbf{Author}: \href{https://sites.google.com/view/fabiofurini/home-page}{\underline{Fabio Furini}}}
\fancyfoot[CO,RE]{\href{https://www.diag.uniroma1.it/}{\includegraphics[scale=0.3]{\logosapienza}}}
% anche le pagine di apertura dei capitoli usano lo stesso piè di pagina
\fancypagestyle{plain}{%
  \fancyhf{}%
  \renewcommand{\headrulewidth}{0pt}%
  \renewcommand{\footrulewidth}{0.4pt}%
  \fancyfoot[LE,RO]{\thepage}%
  \fancyfoot[LO,CE]{\textbf{Author}: \href{https://sites.google.com/view/fabiofurini/home-page}{\underline{Fabio Furini}}}%
  \fancyfoot[CO,RE]{\href{https://www.diag.uniroma1.it/}{\includegraphics[scale=0.3]{\logosapienza}}}%
}


% ---------- palette istituzionale della dispensa ----------
\definecolor{blunotte}{HTML}{16324A}
\definecolor{teal}{HTML}{0E7490}
\definecolor{tealchiaro}{HTML}{E3F2F5}
\definecolor{grigiochiaro}{HTML}{F4F6F7}
\definecolor{verde}{HTML}{1E8449}
\definecolor{verdechiaro}{HTML}{EAF7EE}
\definecolor{giallo}{HTML}{B9770E}
\definecolor{giallochiaro}{HTML}{FEF9E7}
\definecolor{rossomattone}{HTML}{C0392B}
\definecolor{arancio}{HTML}{CA6F1E}
\definecolor{grigiomedio}{HTML}{7F8C8D}
\definecolor{viola}{HTML}{8E44AD}
\definecolor{ocra}{HTML}{B7950B}
\definecolor{codicebg}{HTML}{FBFCFC}

% ---------- titoli ----------
\titleformat{\chapter}[display]
  {\normalfont\huge\bfseries\color{blunotte}}
  {\filleft\Large Chapter \thechapter}{1ex}{\titlerule\vspace{1ex}\filright}
[\vspace{1ex}\titlerule]
\titleformat{\section}{\Large\bfseries\color{teal}}{\thesection}{0.8em}{}
\titleformat{\subsection}{\large\bfseries\color{blunotte}}{\thesubsection}{0.8em}{}

% ---------- box didattici ----------
\newtcolorbox{modello}[1][Model]{%
  enhanced, breakable, colback=green!7!white, colframe=green!55!black,
  title={#1}, fonttitle=\bfseries, boxrule=0.8pt, arc=1.5mm,
  left=1mm, right=1mm, top=1.5mm, bottom=1.5mm}

\newtcolorbox{spiegazione}[1][Line-by-line explanation]{%
  enhanced, breakable, colback=grigiochiaro, colframe=gray!60,
  title={#1}, fonttitle=\bfseries, coltitle=blunotte,
  boxrule=0.5pt, arc=1.5mm, left=2mm, right=2mm}

\newtcolorbox{esempio}[1][Worked example]{%
  enhanced, breakable, colback=teal!6!white, colframe=teal,
  title={#1}, fonttitle=\bfseries, boxrule=0.8pt, arc=1.5mm,
  left=2mm, right=2mm}

\newtcolorbox{insight}[1][Managerial insight]{%
  enhanced, breakable, colback=blunotte!4!white, colframe=blunotte,
  borderline west={2.5pt}{0pt}{blunotte}, sharp corners,
  title={#1}, fonttitle=\bfseries, coltitle=blunotte, boxrule=0pt, frame hidden,
  left=3mm, right=2mm}

\newtcolorbox{attenzione}[1][Warning]{%
  enhanced, breakable, colback=giallochiaro, colframe=giallo,
  borderline west={2.5pt}{0pt}{giallo}, sharp corners,
  title={#1}, fonttitle=\bfseries, boxrule=0pt, frame hidden,
  left=3mm, right=2mm}

\newcounter{esercizio}[chapter]
\renewcommand{\theesercizio}{\thechapter.\arabic{esercizio}}
\newtcolorbox{esercizio}[1][]{%
  enhanced, breakable, colback=white, colframe=blunotte,
  boxrule=0.7pt, arc=1.5mm, left=2mm, right=2mm,
  title={Exercise \refstepcounter{esercizio}\theesercizio\ifstrempty{#1}{}{ --- #1}},
  fonttitle=\bfseries}

% ---------- codice Python ----------
\lstdefinestyle{python}{
  language=Python,
  basicstyle=\ttfamily\small,
  keywordstyle=\color{blunotte}\bfseries,
  stringstyle=\color{verde},
  commentstyle=\color{gray}\itshape,
  numbers=left, numberstyle=\tiny\color{gray}, numbersep=7pt,
  backgroundcolor=\color{codicebg},
  frame=single, rulecolor=\color{gray!40}, framesep=3pt,
  breaklines=true, showstringspaces=false,
  literate={à}{{\`a}}1 {è}{{\`e}}1 {é}{{\'e}}1 {ì}{{\`\i}}1 {ò}{{\`o}}1 {ù}{{\`u}}1
           {À}{{\`A}}1 {È}{{\`E}}1 {²}{{$^2$}}1 {≤}{{$\leq$}}1 {≥}{{$\geq$}}1
           {α}{{$\alpha$}}1 {γ}{{$\gamma$}}1 {λ}{{$\lambda$}}1 {μ}{{$\mu$}}1
           {σ}{{$\sigma$}}1 {ξ}{{$\xi$}}1 {η}{{$\eta$}}1 {π}{{$\pi$}}1 {τ}{{$\tau$}}1
           {ν}{{$\nu$}}1 {ε}{{$\varepsilon$}}1 {ρ}{{$\rho$}}1 {→}{{$\to$}}1 {ombra}{{ombra}}5}
\lstdefinestyle{output}{
  basicstyle=\ttfamily\footnotesize,
  backgroundcolor=\color{grigiochiaro},
  frame=single, rulecolor=\color{gray!40},
  breaklines=true, showstringspaces=false,
  literate={à}{{\`a}}1 {è}{{\`e}}1 {é}{{\'e}}1 {ì}{{\`\i}}1 {ò}{{\`o}}1 {ù}{{\`u}}1 {€}{{\officialeuro}}1}

% ---------- stile comune delle figure pgfplots ----------
% Tutte le figure leggono i CSV in figure/dat/ generati dagli script Python:
% sono quindi completamente modificabili qui in LaTeX (colori, assi, etichette)
% e si aggiornano da sole rieseguendo python3 ../python/esegui_tutti.py
\pgfplotsset{
  lab/.style={
    width=0.74\textwidth, height=6.0cm, grid=major,
    grid style={black!10}, axis line style={black!45},
    tick label style={font=\scriptsize}, label style={font=\small},
    title style={font=\small\bfseries, color=blunotte, align=center},
    legend style={font=\scriptsize, draw=none, fill=none,
      at={(0.5,-0.24)}, anchor=north, /tikz/every even column/.append style={column sep=0.35cm}},
    legend columns=3,
    legend cell align=left,
    cycle list={teal, rossomattone, verde, arancio, blunotte, grigiomedio, viola, ocra},
  },
  labmezza/.style={lab, width=0.495\textwidth, height=5.4cm},
}
\newcommand{\figdat}[1]{figure/dat/#1.csv}
% filtro per riga: scarta i punti la cui colonna #1 non vale #2
\pgfplotsset{
  discard if not/.style 2 args={
    x filter/.append code={%
      \edef\tempa{\thisrow{#1}}\edef\tempb{#2}%
      \ifx\tempa\tempb\else\def\pgfmathresult{nan}\fi
    }
  }
}

% ---------- comandi matematici ricorrenti ----------
\newcommand{\R}{\mathbb{R}}
\newcommand{\Q}{\mathbb{Q}}
\newcommand{\Z}{\mathbb{Z}}
\newcommand{\E}{\mathbb{E}}
\newcommand{\Prob}{\mathbb{P}}
\newcommand{\var}{\mathrm{VaR}}
\newcommand{\cvar}{\mathrm{CVaR}}
\DeclareMathOperator*{\argmin}{arg\,min}
\DeclareMathOperator*{\argmax}{arg\,max}
\usepackage{eurosym}


% ---------- additions of the MIP modelling course ----------
% pseudocode of the heuristics (same package and look as the original exercises)
\usepackage[lined,ruled,linesnumbered]{algorithm2e}
\SetKwInput{Input}{Input}
\SetKwInput{Output}{Output}
\SetKwProg{Function}{Function}{}{}
\definecolor{rossochiaro}{HTML}{FBECEA}
% red box of the algorithms: pseudocode of the constructive heuristics
\newtcolorbox{algoritmo}[1][Algorithm]{%
  enhanced, breakable, colback=red!7!white, colframe=red!55!black,
  title={#1}, fonttitle=\bfseries, boxrule=0.8pt, arc=1.5mm,
  left=1mm, right=1mm, top=1.5mm, bottom=1.5mm}
% statement of a problem (green like the model, title «Problem»)
\newtcolorbox{problema}[1][Problem]{%
  enhanced, breakable, colback=green!7!white, colframe=green!55!black,
  title={#1}, fonttitle=\bfseries, boxrule=0.8pt, arc=1.5mm,
  left=1mm, right=1mm, top=1.5mm, bottom=1.5mm}
% additional modelling questions with solution (white, midnight-blue frame, like the exercises)
\newcounter{domanda}[chapter]
\renewcommand{\thedomanda}{\thechapter.\arabic{domanda}}
\newtcolorbox{domanda}[1][]{%
  enhanced, breakable, colback=white, colframe=blunotte,
  boxrule=0.7pt, arc=1.5mm, left=2mm, right=2mm,
  title={Modelling question \refstepcounter{domanda}\thedomanda\ifstrempty{#1}{}{ --- #1}},
  fonttitle=\bfseries}
\newtcolorbox{soluzione}[1][Solution]{%
  enhanced, breakable, colback=teal!6!white, colframe=teal,
  title={#1}, fonttitle=\bfseries, boxrule=0.8pt, arc=1.5mm,
  left=2mm, right=2mm}
% boolean operators (from the lecture notes on boolean expressions)
\newcommand{\AND}{\mathrel{\texttt{ AND }}}
\newcommand{\OR}{\mathrel{\texttt{ OR }}}
\newcommand{\NOT}{\mathop{\texttt{NOT}}\,}
\newcommand{\true}{\texttt{TRUE}}
\newcommand{\false}{\texttt{FALSE}}
% optimal values and bounds
\newcommand{\zmilp}{z(\mathit{MILP})}
\newcommand{\zlp}{z(\mathit{LP})}
\newcommand{\zdual}{z(\mathit{D}(\mathit{LP}))}
\newcommand{\zlpp}{z(\mathit{LP}^+)}
\newcommand{\zlppp}{z(\mathit{LP}^{++})}
\newcommand{\ub}{\mathit{UB}}
\newcommand{\lb}{\mathit{LB}}
\usepackage{mathrsfs}   % \mathscr{J}_c for the job classes
% instance models with spaced coefficients: shrunk if wider than the text.
% They live in a violet box, different from the green of the general models: it
% tells "the model" from "the model of this instance" at a glance.
\usepackage{adjustbox}
\definecolor{violachiaro}{HTML}{F5EEFA}
\newtcolorbox{istanzabox}[1][The model of the instance]{%
  enhanced, breakable, colback=violachiaro, colframe=viola,
  title={#1}, fonttitle=\bfseries, boxrule=0.8pt, arc=1.5mm,
  left=1mm, right=1mm, top=1.5mm, bottom=1.5mm}
% Le istanze hanno molte colonne: si stringe lo spazio fra colonne invece
% di rimpicciolire il corpo del testo.
\newenvironment{istanza}[1][The model of the instance]
  {\begin{istanzabox}[#1]\footnotesize\setlength{\arraycolsep}{0.7pt}\renewcommand{\arraystretch}{1.4}%
   \begin{center}$\displaystyle}
  {$\end{center}\end{istanzabox}}

% ---------- breathing space of the page ----------
% The pages must not read dense: a touch more leading, a gap between
% paragraphs and a margin around the teaching boxes, which are dense by
% their own nature.
\linespread{1.06}
\setlength{\parskip}{0.35ex plus 0.25ex}
\tcbset{before skip=2.8mm, after skip=2.8mm}
\titlespacing*{\section}{0pt}{2.8ex plus 1ex minus .2ex}{1.4ex plus .2ex}
\titlespacing*{\subsection}{0pt}{2.2ex plus 1ex minus .2ex}{1.1ex plus .2ex}

% ---------- indice analitico ----------
\usepackage{imakeidx}
\makeindex[intoc=true, title=Index, columns=2]
